\section{Approximation Schemes: LO, MF, and SCBA}

\subsection{Why multiple approximation schemes are unavoidable}
The exact interacting central-region problem is defined by the coupled equations
\begin{align}
    G^r &= \sqbrak{(G_0^r)^{-1} - \Sigma_n^r[G]}^{-1},
    \label{eq:exact_dyson_closure_repeat} \\
    G^< &= G^r\paren{\Sigma_L^< + \Sigma_R^< + \Sigma_n^<[G]}G^a.
    \label{eq:exact_keldysh_closure_repeat}
\end{align}
Equation~\eqref{eq:exact_dyson_closure_repeat} is exact, but it is not yet useful computationally because the nonlinear self-energy depends on the very Green's function we are trying to solve for. The same statement applies to Eq.~\eqref{eq:exact_keldysh_closure_repeat}. Therefore the central difficulty of anharmonic NEGF is not writing Dyson and Keldysh equations, but closing them.

To make that precise, let us formally expand the nonlinear self-energy in diagram classes. Symbolically,
\begin{equation}
    \Sigma_n[G]
    =
    \Sigma_{\rm cubic}[G] + \Sigma_{\rm quartic}[G] + \cdots.
    \label{eq:self_energy_split_by_interaction_order}
\end{equation}
Each contribution on the right-hand side is itself an infinite sum of diagrams,
\begin{equation}
    \Sigma_{\rm cubic}[G] = \Sigma_{\rm cubic}^{(2v)}[G] + \Sigma_{\rm cubic}^{(4v)}[G] + \cdots,
\end{equation}
\begin{equation}
    \Sigma_{\rm quartic}[G] = \Sigma_{\rm quartic}^{(1v)}[G] + \Sigma_{\rm quartic}^{(2v)}[G] + \cdots,
\end{equation}
where the superscript indicates the number of interaction vertices retained in the diagrammatic expansion. Thus an approximation scheme must answer two questions.
\begin{enumerate}[label=\arabic*.]
\item Which diagram topologies are kept?
\item When those diagrams are evaluated, are the internal lines harmonic propagators $G_0$, fully dressed propagators $G$, or some effective partially renormalized propagators?
\end{enumerate}
Lowest order, mean field, and SCBA are different answers to those two questions.

\subsection{Lowest order (LO)}
\subsubsection{Definition}
Lowest order means: choose the first nonzero self-energy diagrams generated by the chosen interaction and evaluate all internal lines with the harmonic propagators $G_0$.

Let the first nonvanishing nonlinear self-energy be denoted by $\Sigma_n^{[p]}$, where $p$ is the interaction order in the perturbation series. Then LO means
\begin{equation}
    \Sigma_n^{\rm LO} = \Sigma_n^{[p]}[G_0].
    \label{eq:LO_definition_precise}
\end{equation}
For cubic interactions one typically has $p=2$ because the leading dynamic self-energy comes from two cubic vertices. For quartic Hartree mean field, the first nonzero contribution already appears at one quartic vertex.

To derive how LO enters the Green's function, start from the exact Dyson equation in integral form,
\begin{equation}
    G^r = G_0^r + G_0^r\Sigma_n^r G^r.
    \label{eq:dyson_integral_for_LO_derivation}
\end{equation}
Substitute the right-hand side recursively for $G^r$ itself:
\begin{align}
    G^r
    &= G_0^r + G_0^r\Sigma_n^r\paren{G_0^r + G_0^r\Sigma_n^r G^r} \\
    &= G_0^r + G_0^r\Sigma_n^r G_0^r + G_0^r\Sigma_n^r G_0^r\Sigma_n^r G^r.
\end{align}
Continuing this recursion generates the Neumann series
\begin{equation}
    G^r = G_0^r + G_0^r\Sigma_n^r G_0^r + G_0^r\Sigma_n^r G_0^r\Sigma_n^r G_0^r + \cdots.
    \label{eq:dyson_neumann_series_again}
\end{equation}
If $\Sigma_n^r$ is already small in the perturbative sense, the first correction to $G_0^r$ is the term $G_0^r\Sigma_n^r G_0^r$. Thus the LO retarded Green's function is
\begin{equation}
    G^{r,\rm LO}
    \approx
    G_0^r + G_0^r\Sigma_n^{r,\rm LO}[G_0]G_0^r.
    \label{eq:LO_retarded_from_neumann}
\end{equation}
We can derive the corresponding lesser formula in exactly the same spirit. Start from the exact Keldysh equation,
\begin{equation}
    G^< = G^r\Sigma_{\rm tot}^< G^a,
    \qquad
    \Sigma_{\rm tot}^< = \Sigma_L^< + \Sigma_R^< + \Sigma_n^<.
    \label{eq:exact_keldysh_for_LO_derivation}
\end{equation}
Write
\begin{equation}
    G^r = G_0^r + \delta G^r,
    \qquad
    G^a = G_0^a + \delta G^a,
    \qquad
    \Sigma_{\rm tot}^< = \Sigma_0^< + \Sigma_n^<,
\end{equation}
with
\begin{equation}
    \Sigma_0^< = \Sigma_L^< + \Sigma_R^<.
\end{equation}
Retaining only terms linear in the nonlinear self-energy gives
\begin{align}
    G^<
    &\approx
    G_0^r\Sigma_0^<G_0^a
    + \delta G^r\Sigma_0^<G_0^a
    + G_0^r\Sigma_0^<\delta G^a
    + G_0^r\Sigma_n^<G_0^a.
\end{align}
Now use
\begin{equation}
    G_0^< = G_0^r\Sigma_0^<G_0^a,
\end{equation}
plus
\begin{equation}
    \delta G^r = G_0^r\Sigma_n^{r,\rm LO}G_0^r,
    \qquad
    \delta G^a = G_0^a\Sigma_n^{a,\rm LO}G_0^a.
\end{equation}
Then
\begin{align}
    G^{<,\rm LO}
    &\approx
    G_0^<
    + G_0^r\Sigma_n^{r,\rm LO}G_0^r\Sigma_0^<G_0^a
    + G_0^r\Sigma_0^<G_0^a\Sigma_n^{a,\rm LO}G_0^a
    + G_0^r\Sigma_n^{<,\rm LO}G_0^a \\
    &=
    G_0^<
    + G_0^r\Sigma_n^{r,\rm LO}G_0^<
    + G_0^<\Sigma_n^{a,\rm LO}G_0^a
    + G_0^r\Sigma_n^{<,\rm LO}G_0^a.
    \label{eq:LO_lesser_full_derivation}
\end{align}
Equation~\eqref{eq:LO_lesser_full_derivation} is the explicit linearized Keldysh expansion behind a lowest-order inelastic conductance formula.

\subsubsection{Assumptions}
The assumptions of LO are now transparent from the derivation.
\begin{enumerate}[label=\arabic*.]
\item The retained self-energy $\Sigma_n^{\rm LO}$ is small enough that terms quadratic or higher in $\Sigma_n$ can be dropped in Eq.~\eqref{eq:dyson_neumann_series_again}.
\item The internal lines inside the self-energy diagrams are still well approximated by harmonic propagators.
\item Observables can be trusted even though the feedback of scattering on the internal lines is neglected.
\item The omitted higher-topology diagrams are genuinely subleading in the regime of interest.
\end{enumerate}
Because LO is a truncation before self-consistency, it is generally not a conserving approximation in the strict Baym--Kadanoff sense. It is nevertheless the clearest perturbative benchmark.

\subsubsection{Strengths and limitations}
The strengths of LO follow directly from Eq.~\eqref{eq:LO_retarded_from_neumann} and Eq.~\eqref{eq:LO_lesser_full_derivation}.
\begin{itemize}
\item Every retained correction can be traced to a definite diagram and a definite interaction order.
\item The approximation is analytically transparent and often admits closed-form toy-model benchmarks.
\item Weak-coupling lifetimes derived from LO connect naturally to Fermi's golden rule.
\end{itemize}
The limitations also follow directly from the same equations.
\begin{itemize}
\item Once $\Sigma_n$ is not small, the neglected quadratic and higher terms in Eq.~\eqref{eq:dyson_neumann_series_again} matter.
\item The harmonic internal lines may become a poor representation of the actual broadened spectrum.
\item Conservation relations can be degraded because the self-energy and Green's functions are not solved self-consistently.
\end{itemize}

\subsubsection{Examples from the local paper set}
The local paper set gives three canonical realizations of LO.
\begin{enumerate}[label=\arabic*.]
\item Wang 2007 Fig.~5: a cubic onsite toy model in which the first nontrivial cubic graph is evaluated explicitly.
\item Xu \emph{et al.} 2008: the retarded cubic self-energy is evaluated on harmonic modes to extract linewidths and lifetimes.
\item Wang 2006: the nonlinear correction enters perturbatively in the center region of a mesoscopic junction calculation.
\end{enumerate}
In all three cases the approximation logic is the same: identify the first nonzero diagram class and evaluate it with the harmonic propagators.

\subsection{Mean field (MF)}
\subsubsection{Definition}
Mean field replaces higher-order operator products by lower-order expectation values, thereby reducing the interacting Hamiltonian to an effective quadratic Hamiltonian. The formal pattern is
\begin{equation}
    u_{i_1}\cdots u_{i_m}
    \longrightarrow
    \sum \braket{u\cdots u}\times \text{lower-order operator products}.
    \label{eq:MF_operator_pattern}
\end{equation}
The resulting effective Hamiltonian is harmonic but depends on correlators that must be determined self-consistently.

The cleanest derivation is the quartic case. Start from
\begin{equation}
    H^{(4)} = \frac{1}{4!}\sum_{ijkl}\Phi^{(4)}_{ijkl}u_i u_j u_k u_l.
    \label{eq:quartic_start_for_MF}
\end{equation}
Use the pair decoupling
\begin{align}
    u_i u_j u_k u_l
    \approx
    &\braket{u_i u_j}u_k u_l
    + \braket{u_i u_k}u_j u_l
    + \braket{u_i u_l}u_j u_k \\
    &+ \braket{u_j u_k}u_i u_l
    + \braket{u_j u_l}u_i u_k
    + \braket{u_k u_l}u_i u_j
    - \text{constant terms}.
    \label{eq:quartic_MF_decoupling_repeat}
\end{align}
Substitute Eq.~\eqref{eq:quartic_MF_decoupling_repeat} into Eq.~\eqref{eq:quartic_start_for_MF}. Because the quartic tensor is symmetric, the six contractions produce one common quadratic structure after relabeling dummy indices. Hence
\begin{equation}
    H^{(4)}_{\rm MF}
    =
    \frac{1}{2}\sum_{ij}\Delta K^{(4)}_{ij}u_i u_j + \const.
    \label{eq:quartic_MF_effective_hamiltonian}
\end{equation}
The effective renormalization matrix can be written as
\begin{equation}
    \Delta K^{(4)}_{ij}
    =
    C^{(4)}\sum_{kl}\Phi^{(4)}_{ijkl}\braket{u_k u_l},
    \label{eq:quartic_MF_deltaK_general}
\end{equation}
where $C^{(4)}$ is the convention-dependent combinatorial factor.

Therefore the full effective harmonic matrix is
\begin{equation}
    K_{\rm eff} = K^C + \Delta K^{(4)}.
    \label{eq:effective_harmonic_matrix_MF}
\end{equation}
Once Eq.~\eqref{eq:effective_harmonic_matrix_MF} is known, the effective retarded Green's function is simply
\begin{equation}
    G^r_{\rm MF}(\omega)
    =
    \sqbrak{(\omega+\ii\eta)^2I - K_{\rm eff} - \Sigma_L^r - \Sigma_R^r}^{-1}.
    \label{eq:MF_retarded_green_from_keff}
\end{equation}
The key point is that $K_{\rm eff}$ depends on the equal-time correlator, which depends on the lesser Green's function of the effective problem. Thus mean field is a self-consistent effective harmonic theory.

To connect Eq.~\eqref{eq:quartic_MF_deltaK_general} to NEGF notation, use the identity
\begin{equation}
    G^<_{kl}(t,t) = -\frac{\ii}{\hbar}\braket{u_l(t)u_k(t)}.
\end{equation}
Since equal-time displacement operators commute, we may write
\begin{equation}
    \braket{u_k u_l} = \ii\hbar G^<_{kl}(t,t).
    \label{eq:equal_time_correlation_to_lesser_MF}
\end{equation}
Substituting Eq.~\eqref{eq:equal_time_correlation_to_lesser_MF} into Eq.~\eqref{eq:quartic_MF_deltaK_general} gives a self-energy functional of the form
\begin{equation}
    \Sigma_n^{\rm MF}[G]
    =
    \Sigma_n^{\rm static}[G(t,t)].
    \label{eq:MF_self_energy_functional_form}
\end{equation}
In Wang's quartic tensor convention, Eq.~\eqref{eq:MF_self_energy_functional_form} becomes
\begin{equation}
    \Sigma_{n,jj'}(\tau,\tau') = 3\ii\hbar\,\delta(\tau,\tau')\sum_{kl}T_{jj'kl}G_{kl}(\tau,\tau).
    \label{eq:wang_quartic_MF_closure_repeat}
\end{equation}
Equation~\eqref{eq:wang_quartic_MF_closure_repeat} is therefore not an ad hoc closure; it is the quartic Hartree contraction written in the Wang tensor convention.

\subsubsection{Assumptions}
The assumptions of MF follow directly from Eq.~\eqref{eq:quartic_MF_effective_hamiltonian} and Eq.~\eqref{eq:MF_self_energy_functional_form}.
\begin{enumerate}[label=\arabic*.]
\item Static renormalization is assumed to dominate over dynamic broadening.
\item Higher-order operator products are assumed to be well approximated by pair correlators.
\item The main effect of the interaction is to renormalize the effective force constants or frequencies.
\item The omitted dynamic scattering channels are assumed to be less important than the retained static contractions.
\end{enumerate}
Mean field is therefore most natural when the interaction mainly shifts the spectrum rather than producing strong inelastic linewidths.

\subsubsection{Three-phonon MF}
Cubic interactions can also be decoupled in a mean-field-like way, but the resulting structure is less natural. Start from
\begin{equation}
    H^{(3)} = \frac{1}{3!}\sum_{ijk}\Phi^{(3)}_{ijk}u_i u_j u_k.
    \label{eq:cubic_start_for_MF}
\end{equation}
A naive decoupling gives
\begin{equation}
    u_i u_j u_k
    \approx
    \braket{u_i u_j}u_k + \braket{u_i u_k}u_j + \braket{u_j u_k}u_i
    + \braket{u_i}u_j u_k + \braket{u_j}u_i u_k + \braket{u_k}u_i u_j.
    \label{eq:cubic_MF_decoupling_explicit}
\end{equation}
If the reference state is expanded around a centered equilibrium point, then
\begin{equation}
    \braket{u_i}=0.
\end{equation}
Thus the cubic mean-field correction reduces mainly to a tadpole-like effective linear term proportional to $\braket{u_i u_j}u_k$. Such a term often indicates a shift of the reference equilibrium configuration rather than a faithful description of transport scattering. This is why cubic interactions are usually benchmarked first through LO or SCBA rather than through MF.

\subsubsection{Four-phonon MF}
Quartic interactions are the natural domain of MF because Eq.~\eqref{eq:quartic_MF_effective_hamiltonian} is already quadratic. The self-consistent algorithm follows directly from the equations.
\begin{enumerate}[label=\arabic*.]
\item Start from an initial guess for the equal-time correlator $\braket{u_k u_l}$.
\item Construct $\Delta K^{(4)}$ from Eq.~\eqref{eq:quartic_MF_deltaK_general}.
\item Form $K_{\rm eff}$ from Eq.~\eqref{eq:effective_harmonic_matrix_MF}.
\item Solve for $G^r_{\rm MF}$ from Eq.~\eqref{eq:MF_retarded_green_from_keff}.
\item Compute $G^<_{\rm MF}$ from the Keldysh equation of the effective harmonic problem.
\item Update the equal-time correlator using Eq.~\eqref{eq:equal_time_correlation_to_lesser_MF}.
\item Repeat until $K_{\rm eff}$ or the correlator converges.
\end{enumerate}
This is precisely the logic of self-consistent mean field (SCMF).

\subsubsection{Strengths and limitations}
The strengths of MF are now explicit.
\begin{itemize}
\item It directly captures quartic Hartree renormalization.
\item It reduces the nonlinear problem to a self-consistent harmonic problem.
\item It is often numerically more stable than a fully dynamic self-consistency.
\end{itemize}
Its limitations are equally explicit.
\begin{itemize}
\item It neglects genuine dynamic scattering and linewidth broadening.
\item It may misrepresent transport when the dominant physics is inelastic relaxation rather than static renormalization.
\item For cubic interactions it is usually not the most natural physical approximation.
\end{itemize}

\subsection{Self-consistent Born approximation (SCBA)}
\subsubsection{Definition}
SCBA keeps the Born-level diagram topology but replaces the harmonic internal propagators by dressed propagators determined self-consistently. The difference from LO is therefore not in topology, but in the choice of internal lines.

Suppose the Born self-energy functional is
\begin{equation}
    \Sigma_n^{\rm Born}[X].
\end{equation}
Then LO uses
\begin{equation}
    \Sigma_n^{\rm LO} = \Sigma_n^{\rm Born}[G_0],
    \label{eq:LO_as_Born_on_G0}
\end{equation}
whereas SCBA uses
\begin{equation}
    \Sigma_n^{\rm SCBA} = \Sigma_n^{\rm Born}[G].
    \label{eq:SCBA_as_Born_on_G}
\end{equation}
For example, in a cubic problem the Born structure is quadratic in the internal propagators. Thus,
\begin{equation}
    \Sigma_{\rm cubic}^{\rm Born}[X] \propto \Phi^{(3)} X X \Phi^{(3)}.
\end{equation}
Replacing $X=G_0$ gives LO, while replacing $X=G$ gives SCBA.

Insert Eq.~\eqref{eq:SCBA_as_Born_on_G} into Dyson's equation:
\begin{equation}
    G^r = \sqbrak{(G_0^r)^{-1} - \Sigma_n^{r,\rm Born}[G]}^{-1}.
    \label{eq:SCBA_retarded_fixed_point}
\end{equation}
Likewise, for the lesser sector,
\begin{equation}
    G^< = G^r\paren{\Sigma_L^< + \Sigma_R^< + \Sigma_n^{<,\rm Born}[G]}G^a.
    \label{eq:SCBA_lesser_fixed_point}
\end{equation}
Equation~\eqref{eq:SCBA_retarded_fixed_point} and Eq.~\eqref{eq:SCBA_lesser_fixed_point} define a nonlinear fixed-point problem.

We can derive the iterative structure explicitly. Start from an initial guess
\begin{equation}
    G^{r,(0)} = G_0^r,
    \qquad
    G^{<,(0)} = G_0^<.
    \label{eq:SCBA_initial_guess}
\end{equation}
At iteration $n$, evaluate
\begin{equation}
    \Sigma_n^{r,(n+1)} = \Sigma_n^{r,\rm Born}[G^{(n)}],
    \qquad
    \Sigma_n^{<,(n+1)} = \Sigma_n^{<,\rm Born}[G^{(n)}].
    \label{eq:SCBA_self_energy_update}
\end{equation}
Then update the Green's functions using
\begin{equation}
    G^{r,(n+1)} = \sqbrak{(G_0^r)^{-1} - \Sigma_n^{r,(n+1)}}^{-1},
    \label{eq:SCBA_retarded_update}
\end{equation}
\begin{equation}
    G^{<,(n+1)} = G^{r,(n+1)}\paren{\Sigma_L^< + \Sigma_R^< + \Sigma_n^{<,(n+1)}}G^{a,(n+1)}.
    \label{eq:SCBA_lesser_update}
\end{equation}
The iteration continues until a convergence criterion is met, for example
\begin{equation}
    \norm{G^{r,(n+1)} - G^{r,(n)}} < \varepsilon,
    \qquad
    \norm{G^{<,(n+1)} - G^{<,(n)}} < \varepsilon.
    \label{eq:SCBA_convergence_criterion}
\end{equation}
This is the explicit fixed-point realization of SCBA.

\subsubsection{Assumptions}
The assumptions of SCBA follow directly from Eq.~\eqref{eq:SCBA_as_Born_on_G}.
\begin{enumerate}[label=\arabic*.]
\item The dominant missing physics in LO is the feedback of scattering on the internal propagators.
\item The retained Born topology is assumed to be more important than omitted higher-topology vertex corrections.
\item The self-consistent fixed point is assumed to exist and to represent the physically relevant branch.
\end{enumerate}
Thus SCBA is not ``exact,'' but it is a systematic resummation of the chosen Born diagram class.

\subsubsection{Why SCBA is attractive}
The attraction of SCBA is easiest to see by comparing it with the LO Dyson series. In LO,
\begin{equation}
    G^{r,\rm LO} \approx G_0^r + G_0^r\Sigma_n^{r,\rm LO}[G_0]G_0^r,
\end{equation}
so the internal lines never feel the broadening produced by the self-energy. In SCBA, by contrast, the updated Green's function is immediately recycled into the next self-energy evaluation through Eq.~\eqref{eq:SCBA_self_energy_update}. Therefore broadening, spectral shifts, and nonequilibrium populations feed back onto the scattering kernel itself. This is often essential once the interaction is no longer a tiny perturbation.

\subsubsection{Why SCBA is difficult}
The difficulty of SCBA is also explicit in the equations. Every iteration requires:
\begin{enumerate}[label=\arabic*.]
\item evaluating the retarded nonlinear self-energy,
\item evaluating the lesser nonlinear self-energy,
\item solving Dyson's equation,
\item solving Keldysh's equation,
\item checking convergence and possibly applying mixing or stabilization.
\end{enumerate}
In realistic phonon problems the self-energy evaluation itself contains frequency convolutions. This is why SCBA is substantially more expensive than LO or MF, and why many codes first implement simpler SCBA-like closures before attempting the full contour-resolved version.

\subsection{Comparison table}
\begin{center}
\begin{longtable}{p{0.14\linewidth}p{0.22\linewidth}p{0.24\linewidth}p{0.28\linewidth}}
\toprule
Scheme & Internal lines & Typical physics captured & Main risk \\
\midrule
LO & $G_0$ & first correction, perturbative linewidth, benchmark clarity & misses feedback and may fail at moderate anharmonicity \\
MF & effective static $G$ from renormalized $K$ & frequency renormalization, static quartic effects & misses dynamic broadening and true inelastic channels \\
SCBA & dressed $G$ & feedback of scattering and renormalization, conserving tendencies & numerically expensive and still diagram-class dependent \\
\bottomrule
\end{longtable}
\end{center}
The table is a short summary, but the equation-level distinctions are now clear:
\begin{enumerate}[label=\arabic*.]
\item LO truncates both topology and internal-line feedback.
\item MF keeps a static contraction class and solves the resulting effective harmonic problem self-consistently.
\item SCBA keeps a dynamic Born topology and restores feedback by dressing its internal lines self-consistently.
\end{enumerate}

\subsection{Practical guidance for our code design}
The derivations above translate directly into software design rules.
\begin{enumerate}[label=\arabic*.]
\item The reusable inputs should be the harmonic matrix $K$, the anharmonic tensors $\Phi^{(3)}, \Phi^{(4)}, \ldots$, and the lead self-energies.
\item LO, MF, and SCBA should remain separate modules because each one corresponds to a different closure of the exact Dyson--Keldysh system.
\item Toy models should remain thin wrappers that provide special sparse tensors and paper-specific conventions.
\item The material-ready path should inherit the exact tensor-level derivations, not the toy-model shortcuts used in figure benchmarks.
\end{enumerate}
This separation is the theory basis for keeping the code approximation-first and tensor-first.
